% !TeX root = ../main.tex
\begin{frame}{多目标优化}
  \begin{SlideCanvas}
    \BodyAt{55}{84}{850}{多目标优化共同考虑力学、部署与血流表现，分析性能间的取舍，并以器件模型进一步评价。}
    \PicAt{55}{165}{422}{149}{assets/xue2021_flow.pdf}
    \CaptionAt{55}{329}{422}{Xue等：支架器件中的三维流线}
    \PointAt{55}{374}{422}{壁面剪切应力（WSS）}
    \DetailAt{55}{413}{422}{关注低时间平均值（TAWSS）区域，以及异常高WSS区域的面积。}
    \PointAt{530}{153}{370}{Xue等，2021}
    \DetailAt{530}{191}{370}{负泊松比、刚度与修正渗透率；优化后用器件CFD评价流动。}
    \PointAt{530}{264}{370}{Gharleghi等，2021}
    \DetailAt{530}{302}{370}{7个几何变量、50个设计；低TAWSS／高WSS面积与径向刚度。}
    \PointAt{530}{375}{370}{Zhang等，2026}
    \DetailAt{530}{413}{370}{神经网络与差分进化逆向设计；以球囊—支架—血管模型评价部署。}
    \Sources{\RefXueB，图10及全文；\RefGharleghi。\\\RefZhang。后两项依据期刊公开摘要与章节预览。}
  \end{SlideCanvas}
  \BICConclusionBand{力学与血流目标相互影响，\BICBandEmph{单项最优与整体折中}可对应不同构型}
  \note{Xue等2021年的修正渗透率是流动阻碍相关代理量，不能与WSS或血栓率直接等同。Gharleghi等的7个变量为杆截面、杆尺寸、杆角度、单胞排列、单胞高度、连接形式和连接排列，目标为低TAWSS面积、高WSS面积及径向刚度，50个设计中单项最优与整体折中不同。Zhang、Geng、Yao的论文在线发表于2025年，卷期Biomaterials Advances178属于2026年；使用参数化花瓣手性构型、神经网络与差分进化，不将其称为密度法拓扑优化；期刊摘要及预览支持其器件部署模型，未据此延伸为生物或临床验证。三篇工作说明我们需要定义额外的轨迹定制目标，不能笼统地声称以往只优化单项性能。}
\end{frame}
